\section{Discussion}\label{sec:discussion}

\paragraph{Shared structural assumptions.}
Three assumptions run through all three case studies and deserve to be stated once, prominently. First, \emph{availability is treated as exogenous}: the models condition on the stockout process but do not model why stockouts happen. In reality stockouts concentrate on high-demand days and best-selling items, so the censoring is informative; \citet{braden1991informational} show how censoring changes what observations teach, and the choice-based literature (\S\ref{sec:related}) treats the demand-substitution side our single-item models ignore. The fresh-retail factor curve sitting below the raw empirical sales ratios (Figure~\ref{fig:factor}) is a visible trace of this endogeneity inside a model that does not represent it explicitly. Second, \emph{Gaussian likelihoods are used throughout as a pragmatic compromise}, on counts, on a binary indicator, and on non-negative clipped sales; each case study reported a measurable symptom (negative predictive draws, probability draws outside $[0,1]$, degraded coverage on zero-sales days). Count and truncated likelihoods are the natural refinements, at some computational cost at scale. Third, \emph{availability itself must come from somewhere} at forecast time: the case studies either assume it (scenario inputs) or take it from recorded data (retrospective evaluation); forecasting availability is its own problem and nothing here solves it.

\paragraph{What a controlled comparison would require.}
This paper characterizes mechanisms and their qualitative failure modes; it does not rank them. A ranking would require a dataset recording \emph{both} a binary or fractional availability signal and a per-observation censoring point on the same series, so that gating, censored likelihoods, and multiplicative factors could be cross-applied; a with-and-without-factor ablation at the $50{,}000$-series scale; rolling-origin evaluation rather than the single splits used here; and stress-day scoring defined independently of any candidate model. Until then, claims of the form "mechanism $X$ beats mechanism $Y$" are not supported by the evidence in this paper, and we have tried not to make them.

\paragraph{When each mechanism shades into the next.}
The taxonomy of Table~\ref{tab:taxonomy} has soft edges. A binary flag is the degenerate case of a fractional availability, and the gate $g_t = a_t$ is what the factor \eqref{eq:factor} becomes when $a_t \in \{0,1\}$, the floor is zero, and the model is a recursion rather than a regression. A recorded inventory level is simultaneously a censoring point (when sales hit it) and an availability signal (when it is zero), so the mechanisms compose: one can gate a recursion on the zeros and censor the likelihood at the recorded stock. The practical reading of the taxonomy is not "pick exactly one row" but "let what is recorded decide which components the likelihood can honestly contain."
